\section{El cambio de paradigma de NPT: leer explícitamente los datos de entrenamiento al predecir}

Ahora el objeto de la attention pasa de los tokens a los datapoints. Un parametric model estándar comprime la experiencia en los parámetros $\theta$ tras el entrenamiento y, al hacer inferencia sobre una muestra, por lo general solo lee esa muestra; NPT, en cambio, hace que los puntos de entrenamiento y los de prueba entren juntos al modelo, de modo que la predicción mantiene una dependencia explícita del conjunto de datos visible en ese momento.

\lecturefigure{slide-10-npt-overview.jpg}{Vista general de NPT en una diapositiva: dataset-as-input, datapoint self-attention y stochastic masking.}{Video oficial de Stanford Online, 00:17:00--00:17:21.}

Esta figura presenta tres componentes que no pueden suprimirse: primero, la entrada no es una sola fila, sino un conjunto de datos; segundo, el modelo aplica self-attention entre rows; tercero, el masking objective determina qué valores son visibles y cuáles deben reconstruirse. Si se conservan solo dos de ellos, cualesquiera que sean, no se obtiene lo que el artículo denomina el NPT predictive mechanism.

\lecturefigure{slide-11-npt-motivation.jpg}{Parametric prediction frente a NPT: en NPT, la fila objetivo lee explícitamente otros datapoints.}{Video oficial de Stanford Online, 00:17:22--00:18:17.}

\begin{importantbox}{Definición precisa de «non-parametric» en esta clase}
Non-parametric no equivale a «sin parámetros». NPT conserva una gran cantidad de parámetros aprendibles de attention, embedding y feed-forward. Lo que se subraya aquí es que la función de predicción depende explícitamente, durante la inference, de los supplied training data, lo que se escribe $p(y_*\mid x_*,\mathcal{D}_{\mathrm{train}})$, en lugar de depender solo de $x_*$ y de unos parámetros fijos $\theta$.
\end{importantbox}

\teachervoice{Alguien en redes sociales llamó a NPT «k-NN 2.0». La expresión sirve para formar una intuición: el modelo busca información en otras rows; pero NPT no se limita a buscar vecinos según una distancia fija, sino que puede aprender un proceso de match--transform--aggregate multicapa, dependiente de la tarea y no lineal.}

\subsection{Dónde quedan los non-parametric models tradicionales}

NPT no inventó la non-parametric prediction. Los Gaussian process, los kernel method y los k-nearest neighbors hacen que la predicción dependa explícitamente de los datos de entrenamiento; trabajos como deep kernel learning y neural process también intentan incorporar el representation learning. La tesis de NPT es más acotada: el Transformer puede funcionar como un aprendiz general de relaciones que entrena de extremo a extremo, de forma conjunta, la recuperación, la comparación y la agregación.

\lecturefigure{slide-12-traditional-nonparametric.jpg}{Non-parametric models tradicionales y sus extensiones profundas: el lugar histórico de NPT.}{Video oficial de Stanford Online, 00:22:00--00:23:39.}

\begin{knowledgebox}{Términos clave: tipos de explicit-data prediction}
\begin{tabularx}{\textwidth}{L{0.20\textwidth}L{0.27\textwidth}X}
\toprule
Método & Cómo usa los datos de entrenamiento & Relación con NPT \\
\midrule
k-NN & Recupera vecinos según una distancia predefinida y vota o promedia & NPT puede aprender la similitud, la agregación y las transformaciones posteriores; no se limita a una distancia fija. \\
Kernel method & Pondera los puntos de entrenamiento mediante kernel similarity & La attention también produce pesos que dependen de los datos, pero con representaciones profundas y composición en varias capas. \\
Gaussian process & El kernel define una distribución sobre funciones y proporciona la incertidumbre de la predicción & NPT no es un GP estándar ni posee automáticamente una Bayesian uncertainty coherente. \\
Deep kernel / neural process & Usan redes neuronales para aprender representaciones o una agregación condicional & Comparten con NPT la motivación de «aprendizaje parametrizado de relaciones + context explícito». \\
\bottomrule
\end{tabularx}
\end{knowledgebox}

\teachervoice{Kossen enfatiza que buscaban un non-parametric mechanism fácil de usar, plug-and-play y adaptable a múltiples escenarios, en lugar de depender de un inference scheme especializado y complejo. Este objetivo es una motivación de diseño; no debe interpretarse como prueba de que NPT ya supera a los métodos clásicos en todos los escenarios.}

\begin{warningbox}{NPT no es un sistema de meta-learning ya terminado en esta clase}
Los experimentos del artículo entrenan principalmente un modelo para un conjunto de datos fijo y suponen que train/test provienen de la misma distribución de tareas. El meta-learning suele exigir transferencia entre varias tareas o conjuntos de datos. NPT podría extenderse a un sistema que tome un nuevo dataset como context, pero esa es una línea futura que comentan los ponentes, no una capacidad que los experimentos de esta clase hayan demostrado.
\end{warningbox}

\teachervoice{En la sesión de Q\&A, los tres ponentes trazaron el límite con claridad: el trabajo actual sigue siendo supervised learning estándar; el meta-learning con múltiples conjuntos de datos, zero-shot o pocos gradient step es solo un framing que vale la pena explorar.}

\subsection{Resumen de la sección}

La novedad de NPT no es carecer de parámetros, sino conservar los datapoints de entrenamiento dentro del grafo de cómputo de la predicción y dejar que la attention aprenda a leerlos. Hereda la non-parametric intuition clásica y, al mismo tiempo, amplía el alcance del mecanismo aprendible mediante representaciones profundas y cómputo de relaciones en varias capas.
